
\subsubsection{2.1 Group Relative Policy Optimization}

GRPO was introduced in DeepSeekMath \citep{Shao2024} as an alternative to PPO for large language model fine-tuning. The core innovation is the elimination of the learned value network: advantages are computed group-internally by subtracting the group mean reward and dividing by the group standard deviation. This reduces memory requirements and removes the need to train a critic, at the cost of requiring reward variance within each group for meaningful gradient computation. DeepSeekMath applied GRPO to mathematical reasoning with binary correctness rewards.

Subsequent work has adapted GRPO to code generation. DRIVE \citep{Zhu2025} proposes a curriculum-based difficulty ordering for GRPO code training. GHPO \citep{Guo2025} addresses sparse reward settings in GRPO with heuristic modifications. Both works accept the reward function as a given and focus on algorithmic or curriculum modifications. Neither derives a formal criterion for problem tractability or proposes a prescreening step. The relationship between within-group reward variance and the GRPO gradient magnitude has not been derived explicitly in prior work.

\subsubsection{2.2 Reinforcement Learning for Code Generation}

The application of RL to code generation is well-established. PPOCoder \citep{Shojaee2023} applies PPO with binary execution reward as the primary training signal. CodeRL \citep{Le2022} adds a critic network to reduce variance in the binary reward setting, addressing the variance problem from the algorithmic side rather than the reward design side. Both approaches are effective when the model can occasionally produce fully-passing solutions, but neither provides gradient when the model passes no tests on any rollout.

The most directly relevant recent result is from Afterburner \citep{Du2025}, which reports that SFT initialization is a necessary precondition for GRPO to produce meaningful learning signal on competitive programming tasks. This finding emerged from empirical observation. The present work independently encounters the same prerequisite and additionally provides the analytical framework that characterizes why it exists: without SFT, the model cannot produce partially-correct solutions, so S\_term = 0 for all problems and no prescreened group can be formed.

The APPS benchmark \citep{Hendrycks2021} is the evaluation environment used here. APPS spans introductory to competition difficulty, with the introductory tier (difficulty = 0) being the focus of this work.

\subsubsection{2.3 Reward Signal Design}

The question of binary versus graded reward signals has been studied empirically in code generation. PRLCoder \citep{Ye2025} applies process reward --- rewarding intermediate reasoning steps --- and reports a +5.1\% improvement on MBPP relative to binary outcome reward. This work demonstrates that finer-grained reward signals improve performance, but it operates in a process supervision framework rather than GRPO and does not analyze within-group variance properties.

The theoretical reward shaping literature \citep{Ng1999} establishes conditions under which reward transformations preserve the optimal policy, but this analysis applies to tabular MDPs rather than the group-relative advantage computation in GRPO. Outcome vs. process reward comparisons \citep{Lightman2024} in mathematical reasoning show consistent benefits for process supervision, but again without analysis of the specific variance mechanisms in GRPO group normalization.

The ratio of passed tests to total tests has appeared as an evaluation metric in code generation [Chen et al., 2021; Hendrycks et al., 2021] but its variance properties under GRPO have not been analyzed analytically prior to this work.

\subsubsection{2.4 Positioning}

This work connects reward design to the GRPO gradient equation through a Binomial variance analysis. From the GRPO literature, it takes the group-relative advantage formulation and derives its variance requirements. From the RL-for-code literature, it takes the APPS benchmark and the execution feedback paradigm, and adds a prescreening infrastructure layer. From the reward design literature, it provides the first theoretical derivation of why ratio rewards are structurally advantaged in GRPO's partial-tractability regime specifically.

